% Part: computability
% Chapter: recursive-functions
% Section: pr-functions

\documentclass[../../../include/open-logic-section]{subfiles}

\begin{document}

% SEGMENT OLP-0213-S01

\olfileid{cmp}{rec}{prf}
\olsection{முதன்மை மீள்வரையறைச் சார்புகள்}

முதன்மை மீள்வரையறையையும் சேர்ப்பையும் பயன்படுத்தி, ஏற்கெனவே உள்ள
சார்புகளிலிருந்து புதிய சார்புகளை எவ்வாறு வரையறுக்கலாம் என்பதை மீண்டும்
பதிவுசெய்வோம்.

\begin{defn}
  \ollabel{defn:primitive-recursion} $f$ ஒரு $k$-இடச் சார்பு
  ($k\ge 1$) என்றும், $g$ ஒரு $(k+2)$-இடச் சார்பு என்றும் கொள்க.
  \emph{$f$ மற்றும் $g$ இலிருந்து முதன்மை மீள்வரையறையால்}
  வரையறுக்கப்படும் சார்பு, பின்வரும் சமன்பாடுகளால் வரையறுக்கப்படும்
  $(k+1)$-இடச் சார்பு~$h$:
  \begin{align*}
    h(x_0,\dots,x_{k-1},0) & =  f(x_0,\dots,x_{k-1}) \\
    h(x_0,\dots,x_{k-1},y+1) & = g(x_0,\dots,x_{k-1}, y, h(x_0,\dots,x_{k-1}, y))
  \end{align*}
\end{defn}

\begin{defn}
  \ollabel{defn:composition} $f$ ஒரு $k$-இடச் சார்பு என்றும்,
  $g_0$, \dots, $g_{k-1}$ ஆகியவை அனைத்தும் $n$-இடச் சார்புகளான
  $k$ சார்புகள் என்றும் கொள்க.
  \emph{$f$ உடன் $g_0$, \dots,~$g_{k-1}$ ஆகியவற்றின் சேர்ப்பால்}
  வரையறுக்கப்படும் சார்பு, பின்வருமாறு வரையறுக்கப்படும் $n$-இடச்
  சார்பு~$h$:
  \[
  h(x_0, \dots, x_{n-1}) =
  f(g_0(x_0, \dots, x_{n-1}), \dots, g_{k-1}(x_0, \dots, x_{n-1})).
  \]
\end{defn}

$\Succ$ மற்றும் வீழ்த்தல் சார்புகள்
\[
\Proj{n}{i}(x_0,\dots,x_{n-1}) = x_i,
\]
ஆகியவற்றுடன், ஒவ்வொரு இயல் எண் $n$ மற்றும் $i < n$ க்கும்
சார்பு~$\Zero(x) = 0$ ஐயும் முதன்மை மீள்வரையறைச் சார்புகளில்
சேர்த்துக்கொள்வோம்.

\begin{defn}
  முதன்மை மீள்வரையறைச் சார்புகளின் கணம் என்பது
  $\Nat^n$ இலிருந்து $\Nat$ க்குச் செல்லும் சார்புகளின் கணம்;
  அது பின்வரும் கூறுகளால் தொகுத்தறிதல் முறையில் வரையறுக்கப்படுகிறது:
  \begin{enumerate}
  \item $\Zero$ முதன்மை மீள்வரையறைச் சார்பு.
  \item $\Succ$ முதன்மை மீள்வரையறைச் சார்பு.
  \item ஒவ்வொரு வீழ்த்தல் சார்பு $\Proj{n}{i}$ உம் முதன்மை
    மீள்வரையறைச் சார்பு.
  \item $f$ ஒரு $k$-இட முதன்மை மீள்வரையறைச் சார்பு என்றும்,
    $g_0$, \dots,~$g_{k-1}$ ஆகியவை $n$-இட முதன்மை
    மீள்வரையறைச் சார்புகள் என்றும் இருந்தால், $f$ உடன் $g_0$,
    \dots,~$g_{k-1}$ ஆகியவற்றின் சேர்ப்பு ஒரு முதன்மை
    மீள்வரையறைச் சார்பு.
  \item $f$ ஒரு $k$-இட முதன்மை மீள்வரையறைச் சார்பு என்றும்,
    $g$ ஒரு $k+2$-இட முதன்மை மீள்வரையறைச் சார்பு என்றும்
    இருந்தால், $f$ மற்றும் $g$ இலிருந்து முதன்மை மீள்வரையறையால்
    வரையறுக்கப்படும் சார்பு ஒரு முதன்மை மீள்வரையறைச் சார்பு.
\end{enumerate}
\end{defn}

\begin{explain}
இன்னும் சுருக்கமாக, முதன்மை மீள்வரையறைச் சார்புகளின் கணம் என்பது
$\Zero$, $\Succ$, வீழ்த்தல் சார்புகள்~$\Proj{n}{j}$ ஆகியவற்றைக்
கொண்டதும், சேர்ப்பு மற்றும் முதன்மை மீள்வரையறையின் கீழ் மூடியதுமான
மிகச்சிறிய கணம்.

முதன்மை மீள்வரையறைச் சார்புகளின் கணத்தை “படிநிலைகள்” கொண்டு
வரையறுப்பது அதை விவரிப்பதற்கான மற்றொரு வழி. தொடக்கச் சார்புகளான
$\Zero$, $\Succ$, வீழ்த்தல்கள் ஆகியவற்றின் கணத்தை $S_0$ குறிக்கட்டும்.
இவை படிநிலை~$0$ இன் முதன்மை மீள்வரையறைச் சார்புகள். ஒரு படிநிலை
$S_i$ வரையறுக்கப்பட்டதும்,~$S_i$ இல் ஏற்கெனவே உள்ள சார்புகளுக்கு
சேர்ப்பு அல்லது முதன்மை மீள்வரையறையின் ஒற்றை நிகழ்வைப் பயன்படுத்திக்
கிடைக்கும் அனைத்து சார்புகளின் கணம் $S_{i+1}$ ஆகட்டும். அப்போது
\[
S = \bigcup_{i \in \Nat} S_i
\]
என்பது அனைத்து முதன்மை மீள்வரையறைச் சார்புகளின் கணம்.
\end{explain}

$\Add$ ஒரு முதன்மை மீள்வரையறைச் சார்பு என்பதைச் சரிபார்ப்போம்.

\begin{prop}
  கூட்டல் சார்பு $\Add(x,y) = x+y$ ஒரு முதன்மை மீள்வரையறைச் சார்பு.
\end{prop}

\begin{proof}
$f$ மற்றும்~$g$ என்ற இரு சார்புகளைக் கொண்டு
\olref{defn:primitive-recursion} இன் வடிவத்துடன் பொருந்தும்
$\Add$ இன் முதன்மை மீள்வரையறை ஏற்கெனவே நம்மிடம் உள்ளது:
\begin{align*}
  \Add(x_0, 0) & = f(x_0) = x_0\\
  \Add(x_0, y+1) & = g(x_0, y, \Add(x_0, y)) =
  \Succ(\Add(x_0, y))
\end{align*}
எனவே $f$ மற்றும் $g$ ஆகியவையும் முதன்மை மீள்வரையறைச் சார்புகள்
எனில் $\Add$ அவ்வாறே அமையும். $f(x_0) = x_0 =
\Proj{1}{0}(x_0)$; வீழ்த்தல் சார்புகள் முதன்மை மீள்வரையறைச்
சார்புகளாகக் கணக்கிடப்படுவதால், $f$ ஒரு முதன்மை மீள்வரையறைச்
சார்பு. சார்பு $g$ என்பது $g(x_0, y, z)$ எனும் மூவிடச் சார்பு;
அது பின்வருமாறு வரையறுக்கப்படுகிறது:
\[
g(x_0, y, z) = \Succ(z).
\]
$g$ உம் $\Succ$ உம் முற்றிலும் ஒரே சார்பு அல்ல என்பதால், இது
$g$ முதன்மை மீள்வரையறைச் சார்பு என்று இன்னும் நமக்குக் கூறவில்லை:
$\Succ$ ஓரிடச் சார்பு; $g$ மூவிடச் சார்பாக இருக்க வேண்டும்.
ஆனால் சேர்ப்பால் $g$ ஐ “முறைப்படி” பின்வருமாறு வரையறுக்கலாம்:
\[
g(x_0, y, z) = \Succ(\Proj{3}{2}(x_0, y, z))
\]
$\Succ$ மற்றும் $\Proj{3}{2}$ ஆகியவை முதன்மை மீள்வரையறைச்
சார்புகளாகக் கணக்கிடப்படுவதால், $g$ உம் அவ்வாறே அமைகிறது;
ஏனெனில் முதன்மை மீள்வரையறைச் சார்புகளிலிருந்து சேர்ப்பால் அதை
வரையறுக்கலாம்.
\end{proof}

\begin{prop}
  \ollabel{prop:mult-pr}
  பெருக்கல் சார்பு $\Mult(x,y) = x \cdot y$ ஒரு முதன்மை
  மீள்வரையறைச் சார்பு.
\end{prop}

\begin{proof}
  பயிற்சி.
\end{proof}

\begin{prob}
  \olref[cmp][rec][prf]{prop:mult-pr} ஐ நிரூபிக்க: $\Mult$ இன்
  முதன்மை மீள்வரையறையை
  \olref[cmp][rec][prf]{defn:primitive-recursion} கோரும் வடிவத்தில்
  அமைக்கலாம் என்று காட்டியும், அதற்குரிய சார்புகள் $f$ மற்றும்~$g$
  முதன்மை மீள்வரையறைச் சார்புகள் என்று காட்டுக.
\end{prob}

\begin{ex}
முதன்மை மீள்வரையறைக்கான நமது முதல் எடுத்துக்காட்டு இதோ:
\begin{align*}
h(0) & =  1 \\
h(y+1) & =  2 \cdot h(y).
\end{align*}
$k=0$ என்பதால், இச்சார்பு \olref{defn:primitive-recursion} கோரும்
வடிவத்தில் அமைய முடியாது. இவ்வரையறை மாறிலிகள் $1$ மற்றும் $2$
ஆகியவற்றையும் கொண்டுள்ளது. முதல் சிக்கலைத் தவிர்க்க, போலி உள்ளீடு
ஒன்றை அறிமுகப்படுத்தி, சார்பு~$h'$ ஐ வரையறுப்போம்:
\begin{align*}
h'(x_0, 0) & =  f(x_0) = 1 \\
h'(x_0, y+1) & =  g(x_0, y, h'(x_0, y)) = 2 \cdot h'(x_0, y).
\end{align*}
சார்பு $f(x_0) = 1$ ஐ $\Succ$ மற்றும் $\Zero$ இலிருந்து
சேர்ப்பால் வரையறுக்கலாம்: $f(x_0) = \Succ(\Zero(x_0))$.
சார்பு $g$ ஐ $g'(z) = 2 \cdot z$ மற்றும் வீழ்த்தல்களிலிருந்து
சேர்ப்பால் வரையறுக்கலாம்:
\begin{align*}
g(x_0, y, z) & = g'(\Proj{3}{2}(x_0, y, z))
\intertext{மேலும் $g'$ ஐச் சேர்ப்பால் பின்வருமாறு வரையறுக்கலாம்:}
g'(z) & = \Mult(g''(z), \Proj{1}{0}(z))
\intertext{மேலும்}
g''(z) & = \Succ(f(z)),
\end{align*}
இங்கே $f$ மேலுள்ளவாறு உள்ளது: $f(z) = \Succ(\Zero(z))$.
இப்போது~$h'$ நம்மிடம் இருப்பதால், மீண்டும் சேர்ப்பைப் பயன்படுத்தி
$h(y) = h'(\Proj{1}{0}(y),\Proj{1}{0}(y))$ என வைக்கலாம்.
அடிப்படைச் சார்புகளிலிருந்து சேர்ப்புகள் மற்றும் முதன்மை
மீள்வரையறைகளின் ஒரு தொடரைப் பயன்படுத்தி $h$ ஐ வரையறுக்கலாம்
என்பதை இது காட்டுகிறது; எனவே $h$ ஒரு முதன்மை மீள்வரையறைச் சார்பு.
\end{ex}

\end{document}
